Ce mémoire s'est donné pour objectif d'explorer, de manière progressive et empirique, la mise en pratique de l'apprentissage par renforcement face au défi structurel de la rareté du signal de récompense (\textit{reward sparse}), en partant d'environnements standards à récompense dense pour aboutir à un environnement avancé — Rocket League — où l'objectif final demeure temporellement éloigné des actions élémentaires de l'agent (Section \ref{sec:problematic}).

La Phase I a permis de valider une méthodologie rigoureuse de sélection et d'optimisation de modèles à travers trois algorithmes représentatifs des grandes familles du RL profond moderne : DQN, PPO et SAC \parencite{mnihHumanlevelControlDeep2015,schulmanProximalPolicyOptimization2017,haarnojaSoftActorCriticOffPolicy2018}. Ce choix a mis en évidence des nuances structurelles importantes entre ces approches. D'une part, la distinction \textit{on-policy} / \textit{off-policy} conditionne directement la manière dont les données d'expérience peuvent être réutilisées : DQN et SAC, off-policy, s'appuient sur un \textit{replay buffer} et peuvent réexploiter des transitions passées, tandis que PPO, on-policy, doit intégralement renouveler son batch de données après chaque cycle de mise à jour (Section \ref{sec:rollout}), ce qui limite structurellement sa \textit{sample efficiency} malgré ses multiples époques d'optimisation par rollout. D'autre part, l'espace d'actions constitue un facteur discriminant tout aussi déterminant : DQN est intrinsèquement restreint aux espaces discrets en raison de son opérateur max (Section \ref{sec:dqn}), SAC repose nativement sur des espaces continus, tandis que PPO se distingue par sa polyvalence, applicable indifféremment aux deux types d'espaces grâce à sa paramétrisation d'une distribution de probabilités (catégorielle ou gaussienne) plutôt qu'à une recherche explicite du maximum. Cette polyvalence, conjuguée à sa stabilité empirique élevée, explique en grande partie pourquoi PPO s'est imposé comme un standard industriel du domaine et comme l'algorithme finalement retenu pour l'intégralité du curriculum de la Phase II (Section \ref{sec:discussion:continu-vs-discret}).

L'évaluation multi-seeds de la Phase I (Section \ref{sec:eval_finale}) a par ailleurs révélé une limite méthodologique importante de l'optimisation automatique des hyperparamètres : le score obtenu lors d'une campagne Optuna \parencite{akibaOptunaNextgenerationHyperparameter2019} ne constitue pas un prédicteur fiable de la robustesse d'une configuration une fois celle-ci réentraînée sur un budget plus long et sur des seeds différentes. Le cas de DQN illustre cette limite de façon particulièrement nette : la configuration retenue à l'issue du tuning affichait le deuxième meilleur score (275.14), mais sa récompense moyenne en évaluation finale s'est effondrée à 145.19, avec une seule seed sur cinq atteignant le seuil de résolution. Cet écart s'explique vraisemblablement par une configuration agressive (petit \textit{replay buffer}, fréquence de mise à jour maximale) sur-ajustée au budget de tuning, mais aussi par la sensibilité intrinsèque de l'entraînement à l'initialisation des poids du réseau et à la seed d'échantillonnage, qui introduit une variance substantielle d'une exécution à l'autre indépendamment de la qualité apparente des hyperparamètres sélectionnés. Ce constat invite à considérer tout score d'optimisation automatique comme un indicateur de potentiel plutôt que comme une garantie de performance, et souligne l'importance d'une validation multi-seeds systématique avant de généraliser les conclusions d'une campagne de tuning.

Cette même phase a également permis de dégager un compromis pratique entre SAC et PPO. Sur le plan théorique, SAC est généralement considéré comme l'algorithme le plus intéressant du point de vue de la \textit{sample efficiency} : son \textit{replay buffer} et son double réseau critique lui permettent en principe de réexploiter plus efficacement chaque transition collectée que PPO, qui rejette l'intégralité de son batch après chaque cycle de mise à jour (Tableau \ref{tab:algo-comparison}). Les résultats obtenus sur LunarLanderContinuous vont dans le sens de cette hypothèse : à budget d'entraînement identique, SAC obtient la récompense moyenne la plus élevée aussi bien lors du tuning (284.05, Tableau \ref{tab:scores_optuna}) qu'en évaluation finale (233.44 contre 208.67 pour PPO continu, Tableau \ref{tab:eval_finale}). Cet écart doit néanmoins être interprété avec prudence : compte tenu des écarts-types inter-seeds élevés (45 à 52 points) observés sur seulement 5 seeds (Section \ref{sec:limite-eval-5-seeds}), la différence entre SAC et PPO continu reste du même ordre de grandeur que le bruit inter-seeds et ne permet pas d'établir, sur ce seul environnement, une supériorité statistiquement tranchée — elle est tout au plus directionnellement cohérente avec l'avantage théorique attendu. Ce gain potentiel a par ailleurs un coût plus clairement établi : l'architecture de SAC (deux critiques, un acteur, un mécanisme d'ajustement automatique du coefficient d'entropie à chaque step de gradient) le rend structurellement plus coûteux en temps de calcul que PPO ou DQN (Figure \ref{fig:convergence_boxplot}), un facteur déterminant dans des contextes aux ressources limitées. PPO, à l'inverse, ne bénéficie pas de cette même efficacité par échantillon, mais sa polyvalence face aux espaces d'actions discrets et continus, ainsi que sa stabilité, en font un choix plus généralisable — ce qui a directement motivé sa sélection comme algorithme principal pour la Phase II.

La Phase II a confronté ces enseignements à un environnement bien plus exigeant, où l'objectif de marquer un but constitue un signal de récompense intrinsèquement rare (Section \ref{sec:reward_sparse}). Deux familles de solutions complémentaires ont été mobilisées pour répondre à ce défi. Le \textit{reward shaping} agit sur la composition même de la fonction de récompense : il densifie le signal sparse de l'objectif final en y superposant des signaux intermédiaires continus (rapprochement vers la balle, orientation, gestion du boost), sans en principe modifier la politique optimale recherchée lorsqu'il repose sur une formulation \textit{potential-based} \parencite{645528.657613}. Le \textit{curriculum learning}, en revanche, agit sur la progression temporelle de l'entraînement lui-même : plutôt que de densifier un signal figé, il structure l'apprentissage en une séquence de sous-tâches de complexité croissante, chaque palier s'appuyant sur les compétences déjà consolidées par le précédent \parencite{narvekarCurriculumLearningReinforcement2020}. La nuance entre les deux approches est donc moins une question d'efficacité relative que de niveau d'intervention : le reward shaping enrichit le \textit{quoi} (la fonction de récompense), tandis que le curriculum learning organise le \textit{quand} (la séquence des objectifs). Leur combinaison, plutôt que leur opposition, s'est avérée nécessaire dans ce travail : les premiers essais menés avec l'ensemble des signaux introduits simultanément ont montré une convergence plus lente qu'une introduction progressive et structurée par paliers (Section \ref{sec:curriculum}).

Cette démarche combinée n'a cependant pas été exempte de limites. Le \textit{reward shaping} s'est révélé être une arme à double tranchant : plusieurs occurrences de \textit{reward hacking} \parencite{amodeiConcreteProblemsAI2016} ont été observées en \textit{self-play}, où des signaux individuels mal spécifiés ont été détournés par des stratégies collaboratives dégénérées entre les deux agents (Tableau \ref{tab:reward_hacking}). La correction de ces failles — par une reformulation \textit{zero-sum} des récompenses compétitives, ancrée dans le cadre théorique des jeux de Markov à somme nulle \parencite{littmanMarkovGamesFramework1994} — a constitué une étape méthodologique essentielle pour garantir que l'amélioration d'un agent ne puisse se faire qu'au détriment de son adversaire, condition nécessaire à un entraînement compétitif cohérent en \textit{self-play}.

Sur le plan algorithmique, la comparaison entre PPO et SAC sur Rocket League a confirmé, dans un contexte bien plus exigeant que la Phase I, l'avantage de la polyvalence de PPO sur un espace d'actions discret curaté (90 actions), là où SAC — contraint à l'espace continu brut à huit dimensions et pénalisé par des limites d'implémentation de \textit{rlgym-sac} (Section \ref{sec:limites-sac}) — n'a pas dépassé le stade du contrôle de base. L'agent PPO entraîné sur l'ensemble du curriculum a en revanche atteint un niveau de jeu compétitif validé de manière externe, en battant systématiquement des agents de référence de la communauté (Element, Necto) ainsi qu'une majorité de joueurs humains de niveau Or à Champion, tout en développant des comportements émergents non explicitement récompensés (gestion fine du boost, usage systématique du dérapage) — une illustration concrète du potentiel d'une fonction de récompense bien structurée à faire émerger des stratégies non anticipées par le concepteur (Section \ref{sec:ppo_results}).

Au global, ce travail démontre qu'aucun des trois algorithmes étudiés ne constitue une solution universelle : le choix optimal dépend fondamentalement de la nature de l'espace d'actions, du budget de calcul disponible et de la maturité de l'implémentation logicielle, bien plus que d'un score de tuning isolé. Il illustre également que la rareté du signal de récompense, loin d'être un obstacle insurmontable, peut être adressée efficacement par la combinaison structurée du \textit{reward shaping} et du \textit{curriculum learning}, à condition de rester vigilant face aux risques de \textit{reward hacking} qu'elle introduit. Les pistes d'amélioration identifiées — correction de l'implémentation SAC, préentraînement par imitation, extension à l'entraînement multi-agent (Section \ref{sec:future}) — ouvrent la voie à une validation plus rigoureuse et plus généralisable de ces enseignements pour des travaux futurs.